% Generated draft: requires mathematical and source review.
\documentclass[11pt]{article}
\usepackage[margin=25mm]{geometry}
\usepackage{fontspec}
\setmainfont{Latin Modern Roman}
\usepackage{amsmath,amssymb,mathtools,unicode-math}
\setmathfont{Latin Modern Math}
\usepackage{enumitem,needspace,xurl,hyperref}
\setlength{\parindent}{0pt}
\setlength{\parskip}{5pt}
\setlength{\emergencystretch}{3em}
\newcommand{\N}{\mathbb{N}}
\newcommand{\Z}{\mathbb{Z}}
\newcommand{\Q}{\mathbb{Q}}
\newcommand{\R}{\mathbb{R}}
\newcommand{\C}{\mathbb{C}}
\newcommand{\F}{\mathbb{F}}
\newcommand{\A}{\mathbb{A}}
\newcommand{\PP}{\mathbb{P}}
\newcommand{\E}{\mathbb{E}}
\newcommand{\eps}{\varepsilon}
\newcommand{\dd}{\,\mathrm{d}}
\begin{document}
\section*{UnsolvedMath ID: 30004474}
\textbf{Source problem code:} catalogue entry
\subsection{Definitions and mathematical statement}
Let \(p\) be a prime and let \(X\) be a smooth, connected scheme of finite type over \(\overline{\mathbb F}_p\). The supplied statement uses a parameter space \(S_{\bar\rho}\) for semisimple \(\overline{\mathbb Q}_\ell\)-local systems on \(X\) with residual representation \(\bar\rho\), and a subset \(A_{\bar\rho}\subseteq S_{\bar\rho}\) called its arithmetic points. Semisimplicity means being a finite direct sum of simple local systems. The source record does not specify the rank, the residual representation and its coefficient field, the construction of the parameter space or its Zariski topology, or the definition of arithmetic points. It also does not explicitly state a restriction relating \(\ell\) and \(p\).

Write \(F\) for the Frobenius action intended in the supplied statement. Its definition, including the finite-field model of \(X\) and its compatibility with \(\bar\rho\), is not supplied. With these source-dependent objects and meanings understood, the proposed assertion is the following:
\begin{quote}
For every Zariski-closed subset \(Z\subseteq S_{\bar\rho}\), if there exists an integer \(m\geq 1\) such that \(Z\) is stable under \(F^m\), then
\[
\overline{A_{\bar\rho}\cap Z}^{\,Z}=Z.
\]
\end{quote}
Here the closure is taken in the induced Zariski topology on \(Z\). Equivalently, every nonempty Zariski-open subset of \(Z\) must contain a point of \(A_{\bar\rho}\cap Z\); for \(Z=\varnothing\), this condition is vacuous. The meaning of stability under Frobenius is retained from the source wording: the record does not specify whether it requires equality with the image or only containment. Consequently, this is a statement of the supplied mathematical target subject to unresolved definitions and hypotheses, rather than a fully specified conjecture or an asserted theorem.
\subsection{Short English statement}
The proposed assertion asks whether arithmetic local systems are Zariski dense in every closed locus preserved by some positive power of Frobenius. The parameter space, arithmetic points, residual data, and Frobenius action need further specification before the assertion is mathematically precise. The supplied record labels the problem open; that status is not independently verified.
\subsection{Sources}
\begin{itemize}
\item Supplied local UnsolvedMath registry record, permanent ID \(30004474\), title ``Frobenius-Stable Loci and Dense Arithmetic Points,'' source problem code \texttt{catalogue entry}, in the Oberwolfach Reports Open Problems collection. No individual proposer is supplied; the proposed year is recorded as 2020.
\item Supplied report citation: \emph{Arithmetic Geometry (2021)}, Oberwolfach Reports, pp.~1023--1082, DOI \texttt{10.4171/owr/2020/20}. Supplied source URL and literature reference labeled ``original OWR report'': \url{https://doi.org/10.4171/owr/2020/20}. The supplied background describes Oberwolfach Reports as documenting mathematical workshops and their open-problem discussions. The report has not been independently consulted here, and no precise location within the report is supplied for this assertion.
\item Supplied literature assessment, checked on exactly \texttt{2026-08-22}: status ``open.'' This is supplied dated curation, not independently verified. The assessment describes the material as only a fragment asserting known cases of an unspecified conjecture, with the cases and conjecture statement missing, and says it remains open under the curation rule because it cannot be independently triaged. That description is inconsistent with the explicit density assertion present elsewhere in the record. The background requests independent verification of time-sensitive status and cited claims.
\end{itemize}
\end{document}
